\paragraph{拘束行と主成分の行列式}
全ての表現を共通の係数離散付値環 \(\mathcal O\) 上の格子で実現し、
持ち上げの固有系では主成分と指標の直和の格子を使う。
実際、\(\overline{\mathbb Q}_\ell\) 上で主成分の行列環と \(t'\) 個の体の積に値を
取る表現を組み、その対角実現を考える。その跡は \(N_x\) の跡に一致する。
この表現の像が張る代数は半素である。\(yU(g)y=0\) が全ての \(g\) で成り立つと、
主成分の格子は全整行列環を張るので、主成分では \(U(g)\) の代わりに
任意の行列を入れられる。行列単位 \(E_{ba}\) を入れると \(y_{ab}^2=0\) となる。
一変数の成分では \(g=1\) を代入して同じ結論を得る。
他方、\(N_x\) の像の張る代数上の跡形式は付随表現の定理により非退化である。
従って標数 \(0\) の代数閉体上の Brauer--Nesbitt の定理を適用できる。
詳しくは、二つの表現を組にした像が張る代数は、各成分に上の判定を適用して半単純となる。
その Wedderburn 分解の各単純因子で階数一の冪等元の像の跡を比較すると、
二つの表現でそれぞれの単純加群の重複度が等しくなる。
従って \(N_x\) は対角実現と共役であり、求める整格子が得られる。
\leanlinks{\href{\upstream/grh/src/New/Congruences/FirstOrder/TraceRelations.lean\#L402}{跡の一致から主成分と指標の直和の格子へ}}

以下 \(\rho\) は \(\pi\) の捻られた主成分、\(\gamma=\gamma_\pi\)、
\(a_{\mathrm c}=a_{\mathrm{cyc}}^{-1}\) とする。記号 \(v_{\mathcal O}\) は
一様化元の付値が \(1\) となる付値であり、この規約の長さを
\(v_{\mathcal O}(\ell)\) で割ると \(v_\ell\) に対応する長さになる。
\(\mathcal S=\mathcal O[(X_j)_{j\in J}]\)、\(\mathcal A=\mathcal S[G_E]\) とし、
\(g\in G_E\) を \(\mathcal A\) の基底元にも用いる。

全ての固有系 \(x\) を並べた族 \((N_x(g)-1)_x\) の
\(\mathcal O\)-線形包は有限生成である。
これを生成する有限個の \(g_i\) を、添字集合 \(I\) で並べる。
後で使う旧惰性元と検出点もこの列に入れる。
\(z(g)=(z_i(g))_{i\in I}\in\mathcal O^I\) を
\[
 (N_x(g)-1)_x=\sum_{i\in I} z_i(g)(N_x(g_i)-1)_x
\]
を満たすように選ぶ。行 \(b\in\mathcal A^I\) に対して
\(b\cdot\mathbf y=\sum_i b_i(g_i-1)\) と書く。
\(\mathcal I_\mathrm{aug}\) を \(\mathcal A\) の増大イデアルとし、
\[
 c_\gamma=(\zeta_\pi-1)^{-1},\qquad n_\gamma=c_\gamma(\gamma-1),
 \qquad \rho_0=n_\gamma+(1-n_\gamma)(a_{\mathrm c}-1)
\]
と置く。 \(\rho\) は \(\rho_0\) を単位行列に移し、
指標 \(\chi_a\) は \(\rho_0\) を \(\chi_a(a_{\mathrm c})-1\) に移す。

\(N_\mathrm{row}\) を、\(b\cdot\mathbf y\) が次の四つで生成される左イデアルに
入る行の加群とする。
全ての \(N_x\) で零になる元、\(p_\ast\mathcal I_\mathrm{aug}\)、
新素点の拘束 \(\tau-1-X_jH_j(\phi)(\phi-1)\) とその共役逆元側の拘束、
および \(\mathcal I_\mathrm{aug}\rho_0\) である。
行 \(z(g)\) の値は全ての \(N_x\) の下で \(g-1\) と同じになる。
従って次の行は \(N_\mathrm{row}\) に属する。
\[
 \begin{gathered}
 z(gh)-z(g)-g z(h),\qquad z(g_i)-\mathbf e_i,\qquad
 \sum_g\theta_g z(g)\quad
     \left(\sum_g\theta_g (N_x(g)-1)=0\text{ for all }x\right),\\
 p_\ast b\quad (b\in\mathcal A^I),\qquad
 z(\tau)-X_jH_j(\phi)z(\phi).
 \end{gathered}
\]
最初の三種はどの \(N_x\) でも零になる。
最後の行を \(\mathbf y\) に作用させると、全ての \(N_x\) において拘束そのものになる。
さらに \(u_\gamma=c_\gamma z(\gamma)+(1-n_\gamma)z(a_{\mathrm c})\) は
\(\mathbf y\) に作用させると、全ての \(N_x\) で \(\rho_0\) と同じになる。
従って任意の \(g\) に対して \((g-1)u_\gamma\in N_\mathrm{row}\) である。
\leanlinks{\href{\upstream/grh/src/New/Congruences/FirstOrder/RibetMatrices.lean\#L175}{拘束行と余境界を固定する関係}}

\(\pi\) の主成分で \(X=0\) と置く準同型を
\(\varrho_0\colon\mathcal A\to M_{r'}(\mathcal O)\) とする。
主成分の像が全行列環を張るため、\(\varrho_0\) は全射である。
\(K_{\mathcal O}=\operatorname{Frac}(\mathcal O)\) とし、
すべての行 \(b\in N_\mathrm{row}\) で
\(\sum_i\varrho_0(b_i)w_i=0\) となる
\(w_i\in(K_{\mathcal O}/\mathcal O)^{r'}\) を取る。
\(z_w(g)=\sum_i z_i(g)w_i\) と置く。
上の最初の三種の行から、\(z_w\) は \(\rho\) のコサイクルで、
\(z_w(g_i)=w_i\) であり、全ての \(N_x(g)-1\) の間の線形関係を満たす。
拘束を \(X=0\) で読むと \(z_w(\tau)=0\) である。
\((\gamma-1)u_\gamma\) を読むと \(z(\gamma)\) と同じ行になるので
\(z_w(\gamma)=0\) である。実際 \(\rho(\gamma)-1=(\zeta_\pi-1)I\)
であり、\(1-n_\gamma\) は \(\rho\) の下で零である。
従って \(w\mapsto z_w\) は、行を消す加群からこの可除係数の
コサイクル加群への単射である。後述の局所比較を含む、この加群の
正規化した長さの上界を \(b_{\mathrm{main}}\) とする。
\leanlinks{\href{\upstream/grh/src/New/Congruences/FirstOrder/FittingBlock.lean\#L619}{拘束行を消す加群からコサイクルへ}
  \leanlinksep \href{\upstream/grh/src/New/Congruences/FirstOrder/RibetMatrices.lean\#L539}{主成分への特殊化}}

行を消す加群の長さを行列式へ移すため、
\(v_\ell(p_{\mathrm{row}})>b_{\mathrm{main}}\) となる \(p_{\mathrm{row}}\in\mathcal O\) の冪を取る
（例えば一様化元の \(v_{\mathcal O}(\ell)b_{\mathrm{main}}+1\) 以上の冪）。
\(d=\#I\cdot r'\) とし、ブロック行 \((\varrho_0(b_i))_{i\in I}\)
の \(r'\) 本の行の総和を \(\sigma(b)\in\mathcal O^d\) とする。
\(\sigma(N_\mathrm{row})\) は個々の行も含む。実際、行列単位
\(E_{cc}\) を \(\mathcal A\) に持ち上げて行に左から掛ければ
\(\sigma\) は第 \(c\) 行だけを取り出す。そこで
\[
 L=p_{\mathrm{row}}\mathcal O^d+\sigma(N_\mathrm{row})
\]
を考える。これは \(p_{\mathrm{row}}\mathcal O^d\) を含む全階数の行格子である。
Smith 標準形を用いると、\((\mathcal O/p_{\mathrm{row}})^d\) における
\(L\) の消滅部分の長さは \(\operatorname{length}(\mathcal O^d/L)\) に等しい。
この消滅部分は成分ごとに \(p_{\mathrm{row}}\) で割って \(K_{\mathcal O}/\mathcal O\) に
移すと上の加群に入る。従って \(L\) の基底を行とした行列 \(\mathcal M\) は
\(v_\ell(\det\mathcal M)\le b_{\mathrm{main}}\) を満たす。
各行を \(\sigma(b)\) で置き換えても行列は \(p_{\mathrm{row}}\) を法として等しいので、
\(v_\ell(p_{\mathrm{row}})>b_{\mathrm{main}}\) により行列式の付値は変わらない。

行を \(r'\) 本ずつまとめる。第 \(c\) 行にだけ \(1\) が並ぶ
\(r'\) 次行列 \(F_c\) を \(\mathcal A\) に持ち上げ、
対応する \(b\) に掛けて \(c\) について足せば、
特殊化後の第 \(c\) 行が所定の \(\sigma(b)\) であるブロック行が得られる。
こうして \(N_\mathrm{row}\) の行からなる \(I\) 次行列 \(B_0\) を得る。
その主成分での行列式を \(D_x(X)\) と書くと
\begingroup\postdisplaypenalty=10000
\[
 v_\ell(D_\pi(0))=v_\ell(\det\mathcal M)\le b_{\mathrm{main}} .
\]
\leanlinks{\href{\upstream/grh/src/Lib/General/Algebra/BlockFittingBound.lean\#L31}{DVR 上の行格子とブロック行の持ち上げ}
  \leanlinksep \href{\upstream/grh/src/Lib/General/Algebra/BlockFittingBound.lean\#L152}{可除係数から有限商への移行}}
\endgroup
